% !PRINTED 68
\weberpage{68}

\begin{center}
{\large \S.\ 20.}

\vspace{0.5em}
{\large\gesperrt{Determinanten.}}
\end{center}

Wir betrachten jetzt ein System von \(n^2\) beliebigen Grössen, mit denen die rationalen Rechenoperationen ausgeführt werden können. Zu einer einfachen Bezeichnung dieser Grössen wählen wir einen Buchstaben mit einem doppelten Index \(a_i^{(k)}\), worin \(i\) sowohl als \(k\) die Reihe der Ziffern \(1,2,3,\ldots,n\) durchlaufen soll. Zur besseren Uebersicht ordnen wir diese Grössen in ein Quadrat, so dass alle \(a\) mit demselben oberen Index in einer Horizontalreihe, alle \(a\) mit demselben unteren Index in einer Verticalreihe stehen, und bezeichnen dies Quadrat mit \(\Delta\), also:
\begin{equation}
\Delta=
\begin{vmatrix}
a_1^{(1)}&a_2^{(1)}&a_3^{(1)}&\cdots&a_n^{(1)}\\
a_1^{(2)}&a_2^{(2)}&a_3^{(2)}&\cdots&a_n^{(2)}\\
a_1^{(3)}&a_2^{(3)}&a_3^{(3)}&\cdots&a_n^{(3)}\\
\vdots&\vdots&\vdots&&\vdots\\
a_1^{(n)}&a_2^{(n)}&a_3^{(n)}&\cdots&a_n^{(n)}
\end{vmatrix}.
\tag{1}
\end{equation}
Der Kürze halber nennt man die Horizontalreihen \gesperrt{Zeilen}, die Verticalreihen \gesperrt{Colonnen}.

Wir wollen aber unter dem zwischen verticalen Strichen eingeschlossenen Quadrat nicht nur den Complex der Grössen \(a\) verstehen, sondern eine bestimmte arithmetische Verbindung dieser Grössen, die sich ausrechnen lässt, sobald die \(a\) numerisch gegeben sind, und die wir jetzt beschreiben wollen.

Man bilde das Product der in der von links oben nach rechts unten gehenden Diagonale stehenden Glieder:
\begin{equation}
M=a_1^{(1)}a_2^{(2)}a_3^{(3)}\cdots a_n^{(n)};
\tag{2}
\end{equation}
leite daraus \(\Pi(n)\) Producte \(M,M',M'',\ldots\) her, indem man die unteren Indices permutirt, und gebe jedem so entstandenen Product das \gesperrt{positive oder negative Zeichen}, je nachdem die angewandte Permutation zur ersten oder zur zweiten Art gehört, also nach der Bezeichnung des vorigen Paragraphen
\begin{equation}
M'=\pm a_{\alpha_1}^{(1)}a_{\alpha_2}^{(2)}\cdots a_{\alpha_n}^{(n)}.
\tag{3}
\end{equation}

Die Summe aus diesen Producten
\[
M+M'+M''+\cdots=\Sigma M
\]
%CONT>>
